% Part: many-valued-logic
% Chapter: syntax-and-semantics
% Section: connectives

\documentclass[../../../include/open-logic-section]{subfiles}

\begin{document}

\olfileid{mvl}{syn}{con}

\olsection{语言与联结词}

经典命题逻辑以及许多其他逻辑都使用一套\emph{命题常项}和\emph{联结词}。
例如，我们将以下符号作为原始符号：

\begin{enumerate}
\tagitem{prvFalse}{代表假的命题常项~$\lfalse$。}{}
\tagitem{prvTrue}{代表真的命题常项~$\ltrue$。}{}
\item 逻辑联结词：
  \startycommalist
  \iftag{prvNot}{\ycomma $\lnot$（否定）}{}%
  \iftag{prvAnd}{\ycomma $\land$（合取）}{}%
  \iftag{prvOr}{\ycomma $\lor$（析取）}{}%
  \iftag{prvIf}{\ycomma $\lif$（条件句）}{}%
  \iftag{prvIff}{\ycomma $\liff$（双条件句）}{}%
\end{enumerate}


\iftag{defNot,defOr,defAnd,defIf,defIff,defTrue,defFalse,defEx,defAll}{%
除了上述原始符号，我们也使用定义为缩写的符号，例如
\startycommalist
  \iftag{defNot}{\ycomma $\lnot$（否定）}{}%
  \iftag{defAnd}{\ycomma $\land$（合取）}{}%
  \iftag{defOr}{\ycomma $\lor$（析取）}{}%
  \iftag{defIf}{\ycomma $\lif$（条件句）}{}%
  \iftag{defIff}{\ycomma $\liff$（双条件句）}{}%
  \iftag{defFalse}{\ycomma $\lfalse$（假）}{}%
  \iftag{defTrue}{\ycomma $\ltrue$（真）}.}{}

多值逻辑同样使用这些联结词。然而，在同一逻辑中包含某个联结词的不同版本（例如合取）通常是很有用的，这就需要不同的符号来区分这些版本。
有些多值逻辑还包含在经典逻辑中无等价物的联结词。
因此，我们的处理将比通常更具一般性。

\begin{defn}
  一个\emph{命题语言}由一个\emph{联结词}集合 $\Lang L$ 构成。
  每个联结词 $\star$ 都有一个\emph{元数}；元数为~$n$ 的联结词被称为是\emph{$n$元的}。
  元数为~$0$ 的联结词也称为\emph{常项}；
  元数为~$1$ 的联结词称为\emph{一元的}，
  元数为~$2$ 的则称为\emph{二元的}。
\end{defn}

\begin{ex}
  标准命题逻辑语言 $\Lang L_0$ 由以下联结词（及其对应元数）构成：
  $\lfalse$~($0$),
  $\lnot$~($1$),
  $\land$~($2$),
  $\lor$~($2$),
  $\lif$~($2$)。我们考虑的大多数逻辑都使用这个语言。
  基于传统与约定，有些逻辑会对某些联结词使用不同的符号。
  例如，在乘积逻辑中，合取符号通常是 $\odot$ 而非~$\land$。
  有时，添加新的算子会很方便，例如确定性算子 $\triangle$（$1$元）。
\end{ex}

\end{document}